% Generated draft: requires mathematical and source review.
\documentclass[11pt]{article}
\usepackage[margin=25mm]{geometry}
\usepackage{fontspec}
\setmainfont{Latin Modern Roman}
\usepackage{amsmath,amssymb,mathtools,unicode-math}
\setmathfont{Latin Modern Math}
\usepackage{enumitem,needspace,xurl,hyperref}
\setlength{\parindent}{0pt}
\setlength{\parskip}{5pt}
\setlength{\emergencystretch}{3em}
\newcommand{\N}{\mathbb{N}}
\newcommand{\Z}{\mathbb{Z}}
\newcommand{\Q}{\mathbb{Q}}
\newcommand{\R}{\mathbb{R}}
\newcommand{\C}{\mathbb{C}}
\newcommand{\F}{\mathbb{F}}
\newcommand{\A}{\mathbb{A}}
\newcommand{\PP}{\mathbb{P}}
\newcommand{\E}{\mathbb{E}}
\newcommand{\eps}{\varepsilon}
\newcommand{\dd}{\,\mathrm{d}}
\begin{document}
\section*{UnsolvedMath ID: 30002119}
\textbf{Source problem code:} catalogue entry
\subsection{Definitions and mathematical statement}
The supplied question is: ``Do there exist non-isometric planar domains that are isospectral for the Steklov eigenvalue problem?'' The admissible class of domains is not specified, so the following is a formulation template with unresolved hypotheses.

Let \(\mathcal C\) be a specified class of planar domains \(\Omega\subset\R^2\) for which the unweighted Steklov problem
\[
\begin{cases}
\Delta u=0 & \text{in }\Omega,\\
\partial_{\nu}u=\sigma u & \text{on }\partial\Omega
\end{cases}
\]
is well defined and has a discrete spectrum. Here \(\nu\) is the outward unit normal, \(\partial_\nu u\) is the normal derivative, and an eigenfunction \(u\) must be nonzero and belong to the function space appropriate to the chosen boundary regularity and interpretation of the equations. Write \(\operatorname{Spec}_{\mathrm{St}}(\Omega)\) for the multiset of Steklov eigenvalues, including zero and counting multiplicities. Two such domains are Steklov isospectral when these multisets agree.

If ``non-isometric'' means that no Euclidean rigid motion carries one domain onto the other, the question takes the form
\[
\exists\,\Omega_1,\Omega_2\in\mathcal C:\quad
\operatorname{Spec}_{\mathrm{St}}(\Omega_1)
=\operatorname{Spec}_{\mathrm{St}}(\Omega_2)
\quad\text{and}\quad
\nexists\,Q\in O(2),\ b\in\R^2
\text{ such that }Q\Omega_1+b=\Omega_2.
\]
Here \(O(2)=\{Q\in\R^{2\times2}:Q^{\mathsf T}Q=I\}\).

The source does not determine \(\mathcal C\): boundedness, boundary regularity, connectedness, simple connectivity, and any boundary weight are not explicitly addressed. The displayed boundary condition uses the usual unweighted interpretation, which requires confirmation. The intended meaning of isometry also requires confirmation. Consequently, this template must not be read as selecting a smooth, bounded, simply connected, or otherwise restricted version of the problem.
\subsection{Short English statement}
The question asks whether two planar domains can have the same Steklov eigenvalues, with multiplicities, while being non-isometric. The permitted domains and precise boundary formulation require clarification. The supplied record labels the problem open, but its conflicting verification and literature notes do not establish a verified mathematical status.
\subsection{Sources}
\begin{itemize}
\item Supplied local UnsolvedMath registry record, permanent ID \(30002119\), title \emph{Planar Steklov Isospectral Domains}, exact source code \texttt{catalogue entry}. No individual proposer is supplied; the proposed year is recorded as 2012.
\item Supplied report citation: \emph{Geometric Aspects of Spectral Theory (2013)}, Oberwolfach Reports, pp.\ 2013--2076, DOI \texttt{10.4171/owr/2012/33}. \url{https://doi.org/10.4171/owr/2012/33}. The background attributes the collection to Oberwolfach workshop reports and their open-problem discussions.
\item Supplied literature source, labeled ``OWR report'': \url{https://ems.press/content/serial-article-files/46403}.
\item Literature assessment dated exactly \texttt{2026-08-21}: supplied dated curation, not independently verified. It labels the status open but states that the extracted record contains only historical framing of Kac's question, omits the actual hypotheses or claim, and therefore permits no precise literature-status assessment. The background requests independent verification of time-sensitive status and cited claims.
\item The supplied statement-verification field separately says that the official report was checked, confirms the complete statement, and that the statement was retained byte-for-byte with status ``exact''. This is a supplied verification claim; no independent source inspection is asserted here.
\end{itemize}
\end{document}
